\section{Proofs for Section 5}

This appendix provides the proofs omitted from Section \ref{sec:short_run}.

\subsection{Proof of Proposition \ref{prop:short_run_success}}

By definition,
\[
\Pi_H
=
\bigl[B_1(1)-F\bigr]-B_1(0).
\]
Using \eqref{eq:B1},
\[
B_1(1)=v_LP_A m_A^1,
\qquad
B_1(0)=v_LP_A m_A^0.
\]
Hence
\begin{align}
\Pi_H
&=
v_LP_A m_A^1-F-v_LP_A m_A^0
\nonumber\\
&=
v_LP_A(m_A^1-m_A^0)-F,
\label{eq:app5_PiH}
\end{align}
which is \eqref{eq:DeltaSR}.

Therefore,
\[
\Pi_H>0
\quad\Longleftrightarrow\quad
F<v_LP_A(m_A^1-m_A^0),
\]
which proves \eqref{eq:short_run_condition}.

To derive the threshold in terms of $N_A$, substitute
\[
P_A=\frac{N_A(N_A-1)}{2}
\]
into \eqref{eq:short_run_condition}:
\[
\frac{N_A(N_A-1)}{2}>
\frac{F}{v_L(m_A^1-m_A^0)}.
\]
Equivalently,
\[
N_A^2-N_A-\frac{2F}{v_L(m_A^1-m_A^0)}>0.
\]
The positive root of the associated quadratic equation is
\[
\frac{1}{2}
\left(
1+
\sqrt{1+\frac{8F}{v_L(m_A^1-m_A^0)}}
\right).
\]
Hence, treating $N_A$ as continuous, private support obtains if and only if
\[
N_A\ge
N_{SR}^*
:=
\frac{1}{2}
\left(
1+
\sqrt{1+\frac{8F}{v_L(m_A^1-m_A^0)}}
\right),
\]
with the usual ceiling if $N_A$ must be an integer.
\qed

\subsection{Derivation of the Private--Social Wedge}

Recall the definitions
\[
S_A:=v_LP_A(m_A^1-m_A^0),
\]
\[
D_A:=\beta v_LP_A(m_A^1-m_A^0)\bigl[1-\delta(m_A^1+m_A^0)\bigr],
\]
and
\[
S_{AB}(b):=\beta\eta v_XQ_{AB}b-\frac{\kappa}{2}b^2.
\]

Since
\[
\Pi_H=S_A-F
\]
and
\[
\Delta_H
=
v_LP_A(m_A^1-m_A^0)\Bigl[(1+\beta)-\beta\delta(m_A^1+m_A^0)\Bigr]-F,
\]
we obtain
\begin{align}
\Delta_H-\Pi_H
&=
\beta v_LP_A(m_A^1-m_A^0)\bigl[1-\delta(m_A^1+m_A^0)\bigr]
\nonumber\\
&=D_A.
\label{eq:app5_wedge}
\end{align}
Hence
\[
\Delta_H=\Pi_H+D_A,
\]
which is \eqref{eq:dynamic_wedge}.

Likewise, from Section \ref{sec:refresh},
\[
\Delta_{H+b}(b)=\Delta_H+S_{AB}(b),
\]
so
\[
\Delta_{H+b}(b)=\Pi_H+D_A+S_{AB}(b),
\]
which is \eqref{eq:bundle_wedge}.

Finally, using the optimal pipeline policy from Section \ref{sec:refresh},
\[
R_{AB}(N_A,N_B)
=
\frac{(\beta\eta v_XQ_{AB})^2}{2\kappa},
\]
we obtain
\[
\Delta_{H+b}^*=\Pi_H+D_A+R_{AB}(N_A,N_B),
\]
which is \eqref{eq:bundle_wedge_star}.

\subsection{Proof of Proposition \ref{prop:dynamic_inefficiency}}

Recall
\[
\Pi_H=S_A-F,
\qquad
\Delta_H=S_A+D_A-F,
\qquad
\Delta_{H+b}^*=S_A+D_A+R_{AB}(N_A,N_B)-F.
\]
Hence
\[
F_H^*(N_A)=S_A+D_A,
\qquad
F_{H+b}^*(N_A,N_B)=S_A+D_A+R_{AB}(N_A,N_B).
\]

\paragraph{Part 1: Private over-support of a hub alone.}
If
\[
\delta(m_A^1+m_A^0)>1,
\]
then
\[
1-\delta(m_A^1+m_A^0)<0,
\]
so by definition,
\[
D_A=\beta v_LP_A(m_A^1-m_A^0)\bigl[1-\delta(m_A^1+m_A^0)\bigr]<0.
\]
Therefore
\[
\Delta_H=\Pi_H+D_A<\Pi_H.
\]

Next, the hub is privately supported under the short-run rule if and only if
\[
\Pi_H>0
\quad\Longleftrightarrow\quad
F<S_A.
\]
The hub is socially undesirable as a stand-alone policy if and only if
\[
\Delta_H<0
\quad\Longleftrightarrow\quad
F>F_H^*(N_A).
\]
Combining these two conditions gives
\[
\max\{0,F_H^*(N_A)\}<F<S_A.
\]
It remains to show that this interval is nonempty. Since $D_A<0$,
\[
F_H^*(N_A)=S_A+D_A<S_A.
\]
Hence
\[
\max\{0,F_H^*(N_A)\}<S_A,
\]
so the interval is nonempty. This proves \eqref{eq:dynamic_inefficiency_interval}.

\paragraph{Part 2: Private under-support of a hub-plus-pipeline bundle.}
If
\[
D_A+R_{AB}(N_A,N_B)>0,
\]
then
\[
F_{H+b}^*(N_A,N_B)=S_A+D_A+R_{AB}(N_A,N_B)>S_A.
\]

The hub is not privately supported under the short-run rule if and only if
\[
\Pi_H\le 0
\quad\Longleftrightarrow\quad
F\ge S_A.
\]
The hub-plus-pipeline bundle is socially desirable if and only if
\[
\Delta_{H+b}^*>0
\quad\Longleftrightarrow\quad
F<F_{H+b}^*(N_A,N_B).
\]
Combining these gives
\[
S_A\le F<F_{H+b}^*(N_A,N_B),
\]
which is exactly \eqref{eq:under_support_interval}. Because
\[
F_{H+b}^*(N_A,N_B)>S_A,
\]
this interval is nonempty. This completes the proof.
\qed
